\begin{tikzpicture}[
  >=Latex,
  x=0.65cm,
  y=0.65cm,
  axis/.style={->,black!65,line width=0.6pt},
  firstline/.style={THMBlue,line width=1.15pt},
  secondline/.style={olive!80!black,line width=3pt},
  tick/.style={black!65,line width=0.5pt}
]
  \draw[axis] (-2.45,0)--(2.65,0) node[right] {$x_1$};
  \draw[axis] (0,-2.45)--(0,2.65) node[above] {$x_2$};
  \foreach \x in {-2,-1,1,2}{
    \draw[tick] (\x,0.06)--(\x,-0.06);
    \draw[tick] (0.06,\x)--(-0.06,\x);
  }
  % Die breite olivfarbene und die schmale blaue Linie liegen exakt
  % aufeinander, sodass beide identischen Geraden sichtbar bleiben.
  \draw[secondline] (-2.3,2.3)--(2.3,-2.3);
  \draw[firstline] (-2.3,2.3)--(2.3,-2.3);
\end{tikzpicture}
